\clearpage
\SourceLeaf{189}
\pdfbookmark[1]{Einrichtung und Gebrauch der Tafeln.}{witwen-einrichtung}
\begin{center}EINRICHTUNG UND GEBRAUCH DER TAFELN.\end{center}
\noindent\hspace*{1.2em} [Den Zahlenangaben dieser Tafeln liegen die von \textsc{Brune} im 16. Bande des \textsc{Crelle}schen Journals für Mathematik zusammengestellten Erfahrungen über die in der k. Preussischen allgemeinen Witwen-Verpflegungs-Anstalt während der Zeit von 1776 bis 1834 successive aufgenommenen 31500 Ehepaare zu Grunde. Es sind hier angegeben die Logarithmen der Anzahl der Frauen (\(\log fm\)) und der Männer (\(\log FM\)), welche unter 10000, die das vollendete 20\textsuperscript{ste} Lebensjahr erreichten, bis zu dem Ende des in der Mitte bemerkten Altersjahres (\(m\) oder \(M\)) gelangten, jedoch mit der Abweichung von \textsc{Brune}, dass das Absterben der Männer über 80 Jahren nach demselben Verhältnisse gerechnet ist, welches jenen Erfahrungen gemäss bei dem weiblichen Geschlechte gilt; weil wie in der Bilanzrechnung von 1845 erwähnt wird, die Registratur der Preussischen Witwenkasse zur directen Bestimmung des Absterbens der Männer im hohen Alter keine hinreichenden Daten enthält. Neben den Logarithmen der Lebenden stehen unter der Überschrift decr. die absoluten Werthe der Unterschiede jener Logarithmen (\(\log gm=\log\frac{f(m-1)}{fm}\) und \(\log Gm=\log\frac{F(M-1)}{FM}\)) in Einheiten der siebenten Decimale ausgedrückt. Mit Hülfe der so erhaltenen Tafel für die Sterblichkeit hat \textsc{Gauss} die einfachen und die Verbindungsrenten bei dem Zinsfuss von \(3\tfrac{1}{2}\) Proc. und von 4 Proc. berechnet.

Die sowohl in Logarithmen als in Zahlen (\(\varphi m,\Phi M\)) dargestellten einfachen Leibrentenwerthe gelten für das Ende des in der Mitte angegebenen Lebensjahres der Frau (\(m\)) oder des Mannes (\(M\)) als jetzigen Zeitmoment und unter der Voraussetzung, dass für den Fall des Erlebens des Endes jedes der nachfolgenden Jahre dann die Münzeinheit gezahlt wird, so dass also
\[
\begin{aligned}
\varphi m\mathord{.}fm&=\rho f(m+1)+\rho\rho f(m+2)+\rho^3f(m+3)+.\ .\ .\\
\Phi M\mathord{.}FM&=\rho F(M+1)+\rho\rho F(M+2)+\rho^3F(M+3)+.\ .\ .
\end{aligned}
\]
ist, wenn \(\rho\) den Discontofactor (\(=\tfrac{200}{207}\) bei \(3\tfrac{1}{2}\) Proc. und \(=\tfrac{25}{26}\) bei 4 Proc.) bezeichnet.

Die Tafel der Verbindungsrenten enthält die Logarithmen der Werthe \(\psi(m,M)\), welche für den Zeitpunkt des zur Seite stehenden Alters (\(M\)) des Mannes solchen Renten, die am Schlusse jedes der folgenden Jahre im Falle des gleichzeitigen Lebens des Mannes (vom jetzigen Alter \(=M\)) und der Frau (vom Altersunterschiede \(=m-M\)) mit der Münzeinheit gezahlt werden, gleich kommen und also durch die Formel bestimmt sind:
\[
\begin{aligned}
\psi(m,M)\mathord{.}fm\mathord{.}Fm&=\rho f(m+1)\mathord{.}F(M+1)\\
&\quad+\rho\rho f(m+2)\mathord{.}F(M+2)\\
&\quad+\rho^3f(m+3)\mathord{.}F(M+3)+.\ .
\end{aligned}
\]
